\scanpage{427}
\mspage{339}

qui est ``$\alpha_0$-$\pi$-adaptée'' i.e. telle que \deleted{\unsure{2) les $x_i\in$}} \deleted{1) les $x_i\in\mathrm{Ob}\,\unsure{M|\mathcal M}_\pi$ 2)} $1^\circ)$ $I$ \inserted{équiv. à une petite cat. et} \unsure{grd} devant $\pi$, $2^\circ)$ les $\varinjlim$ de type $I_{\alpha_0}$ existent dans $I$, $3^\circ)$ le foncteur $I\longrightarrow M$, $i\mapsto x_i$, commute aux $\varinjlim$ de type $I_{\alpha_0}$, et enfin $4^\circ)$ les $x_i\in\mathrm{Ob}\,M_\pi$, (comparer 4.13.11), on a ceci (où on désigne par $I_{\underline S}$ la sous-catégorie str. pleine de $I$ formée des $i$ tels que $x_i\in\mathrm{Ob}\,\underline S$) : on a
\[
 x\in\mathrm{Ob}\,\underline S\quad\text{\unsure{sss|ssi}}\quad I_{\underline S}\text{ est cofinal }\unsure{\text{dans}}\,I.
\]

c') \inserted{(\unsure{et|{} }\uuline{si} $I_{\alpha_0}$ est grd devant $\pi_0$ et $\pi^{\pi_0}=\pi$, pour pouvoir invoquer 4.13.14), cf. p. 355\unsure{$>$|.}} Comme c), \deleted{mais} en supposant $I$ \unsure{ordonnée|ordonné}.

% slanted
\margin{\deleted{\unsure{Attention|Ajouter|NB}}}

\margin{\uline{NB} Si on ne suppose pas $I_{\alpha_0}$ grd devant $\pi_0$ et $\pi^{\pi_0}=\pi$, alors seules les conditions de a) à c) (a, a', b, c) sont équivalentes, et \uline{impliquent} c', mais ne \uline{\unsure{lui|leur}} sont \unsure{peut être|peut-être} pas équivalentes.}

Démonstration. Elle se fait suivant le $=$ diagramme logique que celle de 4.13.15
\diagram{\begin{tikzcd}[ampersand replacement=\&,column sep=small,row sep=small]
 \& \& \& c'\arrow[dr,Rightarrow,"!"{name=bang}] \& \& {\begin{gathered}\text{ici, il faut supposer}\\ I_{\alpha_0}\text{ grd devant }\pi_0\\ \text{et }\pi^{\pi_0}=\pi\end{gathered}}\\
 a\arrow[r,Rightarrow] \& a'\arrow[r,Rightarrow,"!"] \& c\arrow[ur,Rightarrow]\arrow[dr,Rightarrow] \& \& a'\arrow[r,Rightarrow] \& {a,}\\
 \& \& \& b\arrow[ur,Rightarrow] \& \& {}
 \arrow[from=bang,to=1-6,end anchor=west]
\end{tikzcd}}
où on a marqué du signe ! les deux implications qui sont ``vaches'', en ce sens qu'elles font appel au lemme-clef 4.13.17 (p\unsure{g|ages} 329-330), qui se prouve par une construction transfinie un peu \unsure{vissée|vicieuse}. La démonstration \deleted{\unsure{suit}} des sept implications \inserted{du diagramme} est essentiellement celle donnée dans les pages précédentes 311-337,

% archive sheet number: 426

